% GENERATED FILE. Edit canonical lessons, then run npm run generate:books.
% canonical-source-hash: fnv1a64:5b08f72048478a12
% canonical-lessons: HI-C457-bhay, HI-C457-ghabrana, HI-C457-darana, HI-C457-kayar, HI-C457-siharna

\chapter{Fear}
\label{ch:hi-fear}
\hlchaptermodality{457}

\begin{chapteropening}
\textbf{By the end of this chapter:} \emph{I can say what frightened someone.}

Complete the last lesson of chapter 457: I can say what frightened someone.
\end{chapteropening}

\section[\texorpdfstring{\hi{भय} (bhay)}{भय (bhay)}]{\hi{भय} (bhay) — dread; fear}

\label{lesson:HI-C457-bhay}



\begin{quote}
Before the new one: say the Hindi for to fly into a rage, then the Hindi for to get fed up.
\end{quote}

\subsection*{You'll want to know: \hi{भय}}
\textbf{\hi{भय}} — \emph{bhay} — \textquotedblleft{}dread; fear\textquotedblright{}.

\begin{quote}
\hi{बच्चों} \hi{को} \hi{साँप} \hi{से} \hi{बहुत} \hi{भय} \hi{लगता} \hi{है।}
\end{quote}

\begin{quote}
\emph{baccoṁ ko sā̃p se bahut bhay lagtā hai.}
\end{quote}

\begin{quote}
Children are terribly afraid of snakes.
\end{quote}

Masculine.

\begin{cousinweb}
\textbf{\hi{भय}} (\emph{bhay}) is a tatsama, borrowed from Sanskrit \emph{bhaya}, \textquotedblleft{}fear, alarm, dread\textquotedblright{} (Monier-Williams). It is more formal than \emph{ḍar}.
\end{cousinweb}

\begin{grammarlens}[title={What you've built}]
One more word for describing what you lived and felt.
\end{grammarlens}

\subsection*{Your turn}
\begin{itemize}
\raggedright
  \item \emph{Say it:} \emph{bhay}
  \item \emph{Say it:} \emph{bhay}, once more
  \item \emph{Say it:} \emph{baccoṁ ko sā̃p se bahut bhay lagtā hai.} — Children are terribly afraid of snakes.
  \item \emph{Recall:} say the Hindi for to fly into a rage, then the Hindi for to get fed up, then say \emph{bhay} again
\end{itemize}

\begin{tcolorbox}[breakable,colback=teal!4,colframe=teal!35!black,title={Before you move on}]
What does \hi{भय} mean? (\textquotedblleft{}Dread; fear\textquotedblright{}.) Say it once more.
\end{tcolorbox}
\section[\texorpdfstring{\hi{घबराना} (ghabrānā)}{घबराना (ghabrānā)}]{\hi{घबराना} (ghabrānā) — to panic; to get flustered}

\label{lesson:HI-C457-ghabrana}



\begin{quote}
Before the new one: say the Hindi for to get fed up, then the Hindi for dread.
\end{quote}

\subsection*{You'll want to know: \hi{घबराना}}
\textbf{\hi{घबराना}} — \emph{ghabrānā} — \textquotedblleft{}to panic; to get flustered\textquotedblright{}.

\begin{quote}
\hi{परीक्षा} \hi{से} \hi{पहले} \hi{अरुण} \hi{घबरा} \hi{गया।}
\end{quote}

\begin{quote}
\emph{parīkṣā se pahle Aruṇ ghabrā gayā.}
\end{quote}

\begin{quote}
Arun panicked before the exam.
\end{quote}

Past (he …): \hi{वह} \hi{घबरा} \hi{गया}, \emph{vah ghabrā gayā}.

\begin{cousinweb}
\textbf{\hi{घबराना}} (\emph{ghabrānā}) is probably inherited: CDIAL lists Hindi \emph{ghābrā}, \textquotedblleft{}startled\textquotedblright{}, under \emph{*ghabbhara}, a form it reconstructs (CDIAL 4427), and \emph{ghabrānā} belongs with it. \emph{ghabrāhaṭ}, \textquotedblleft{}nervousness\textquotedblright{}, is its noun.
\end{cousinweb}

\begin{grammarlens}[title={What you've built}]
One more word for describing what you lived and felt.
\end{grammarlens}

\subsection*{Your turn}
\begin{itemize}
\raggedright
  \item \emph{Say it:} \emph{ghabrānā}
  \item \emph{Say it:} \emph{ghabrānā}, once more
  \item \emph{Say it:} \emph{parīkṣā se pahle Aruṇ ghabrā gayā.} — Arun panicked before the exam.
  \item \emph{Recall:} say the Hindi for to get fed up, then the Hindi for dread, then say \emph{ghabrānā} again
\end{itemize}

\begin{tcolorbox}[breakable,colback=teal!4,colframe=teal!35!black,title={Before you move on}]
What does \hi{घबराना} mean? (\textquotedblleft{}To panic; to get flustered\textquotedblright{}.) Say it once more.
\end{tcolorbox}
\section[\texorpdfstring{\hi{डराना} (ḍarānā)}{डराना (ḍarānā)}]{\hi{डराना} (ḍarānā) — to frighten}

\label{lesson:HI-C457-darana}



\begin{quote}
Before the new one: say the Hindi for dread, then the Hindi for to panic.
\end{quote}

\subsection*{You'll want to know: \hi{डराना}}
\textbf{\hi{डराना}} — \emph{ḍarānā} — \textquotedblleft{}to frighten\textquotedblright{}.

\begin{quote}
\hi{अरुण} \hi{ने} \hi{कहानी} \hi{सुनाकर} \hi{छोटे} \hi{भाई} \hi{को} \hi{डराया।}
\end{quote}

\begin{quote}
\emph{Aruṇ ne kahānī sunākar choṭe bhāī ko ḍarāyā.}
\end{quote}

\begin{quote}
Arun frightened his little brother with a story.
\end{quote}

Past, with \hi{ने} (he or she …): \hi{उसने} \hi{भाई} \hi{को} \hi{डराया}, \emph{usne bhāī ko ḍarāyā}.

\begin{cousinweb}
\textbf{\hi{डराना}} (\emph{ḍarānā}) is the causative partner of \emph{ḍarnā}, \textquotedblleft{}to be afraid\textquotedblright{}, which CDIAL lists under Sanskrit \emph{darati}, \textquotedblleft{}fears\textquotedblright{} (CDIAL 6190): to make someone afraid.
\end{cousinweb}

\begin{grammarlens}[title={What you've built}]
One more word for describing what you lived and felt.
\end{grammarlens}

\subsection*{Your turn}
\begin{itemize}
\raggedright
  \item \emph{Say it:} \emph{ḍarānā}
  \item \emph{Say it:} \emph{ḍarānā}, once more
  \item \emph{Say it:} \emph{Aruṇ ne kahānī sunākar choṭe bhāī ko ḍarāyā.} — Arun frightened his little brother with a story.
  \item \emph{Recall:} say the Hindi for dread, then the Hindi for to panic, then say \emph{ḍarānā} again
\end{itemize}

\begin{tcolorbox}[breakable,colback=teal!4,colframe=teal!35!black,title={Before you move on}]
What does \hi{डराना} mean? (\textquotedblleft{}To frighten\textquotedblright{}.) Say it once more.
\end{tcolorbox}
\section[\texorpdfstring{\hi{कायर} (kāyar)}{कायर (kāyar)}]{\hi{कायर} (kāyar) — cowardly; a coward}

\label{lesson:HI-C457-kayar}



\begin{quote}
Before the new one: say the Hindi for to panic, then the Hindi for to frighten.
\end{quote}

\subsection*{You'll want to know: \hi{कायर}}
\textbf{\hi{कायर}} — \emph{kāyar} — \textquotedblleft{}cowardly; a coward\textquotedblright{}.

\begin{quote}
\hi{मेरा} \hi{कुत्ता} \hi{बहुत} \hi{कायर} \hi{है}, \hi{चूहे} \hi{से} \hi{भी} \hi{डरता} \hi{है।}
\end{quote}

\begin{quote}
\emph{merā kuttā bahut kāyar hai, cūhe se bhī ḍartā hai.}
\end{quote}

\begin{quote}
My dog is a real coward; it's even afraid of a mouse.
\end{quote}

\begin{cousinweb}
\textbf{\hi{कायर}} (\emph{kāyar}) is a tadbhava: CDIAL lists Hindi \emph{kāyar}, \textquotedblleft{}coward\textquotedblright{}, under Sanskrit \emph{kātara}, \textquotedblleft{}cowardly, timid\textquotedblright{} (CDIAL 3027); the \emph{t} between vowels wore away to \emph{y}.
\end{cousinweb}

\begin{grammarlens}[title={What you've built}]
One more word for describing what you lived and felt.
\end{grammarlens}

\subsection*{Your turn}
\begin{itemize}
\raggedright
  \item \emph{Say it:} \emph{kāyar}
  \item \emph{Say it:} \emph{kāyar}, once more
  \item \emph{Say it:} \emph{merā kuttā bahut kāyar hai, cūhe se bhī ḍartā hai.} — My dog is a real coward; it's even afraid of a mouse.
  \item \emph{Recall:} say the Hindi for to panic, then the Hindi for to frighten, then say \emph{kāyar} again
\end{itemize}

\begin{tcolorbox}[breakable,colback=teal!4,colframe=teal!35!black,title={Before you move on}]
What does \hi{कायर} mean? (\textquotedblleft{}Cowardly; a coward\textquotedblright{}.) Say it once more.
\end{tcolorbox}
\section[\texorpdfstring{\hi{सिहरना} (siharnā)}{सिहरना (siharnā)}]{\hi{सिहरना} (siharnā) — to shudder; to shiver}

\label{lesson:HI-C457-siharna}



\begin{quote}
Before the new one: say the Hindi for to frighten, then the Hindi for cowardly.
\end{quote}

\subsection*{You'll want to know: \hi{सिहरना}}
\textbf{\hi{सिहरना}} — \emph{siharnā} — \textquotedblleft{}to shudder; to shiver\textquotedblright{}.

\begin{quote}
\hi{साँप} \hi{देखकर} \hi{मीरा} \hi{सिहर} \hi{उठी।}
\end{quote}

\begin{quote}
\emph{sā̃p dekhkar Mīrā sihar uṭhī.}
\end{quote}

\begin{quote}
Mira shuddered when she saw the snake.
\end{quote}

Past (she …): \hi{वह} \hi{सिहर} \hi{उठी}, \emph{vah sihar uṭhī}.

\begin{cousinweb}
\textbf{\hi{सिहरना}} (\emph{siharnā}) is inherited: CDIAL lists Hindi \emph{siharnā}, \textquotedblleft{}to shiver\textquotedblright{}, under Sanskrit \emph{śikhara}, a word for a peak or point, perhaps through the hair standing on end (CDIAL 12435). Fear, cold and delight all make one shudder with it.
\end{cousinweb}

\begin{grammarlens}[title={What you've built}]
One more word for describing what you lived and felt.
\end{grammarlens}

\subsection*{Your turn}
\begin{itemize}
\raggedright
  \item \emph{Say it:} \emph{siharnā}
  \item \emph{Say it:} \emph{siharnā}, once more
  \item \emph{Say it:} \emph{sā̃p dekhkar Mīrā sihar uṭhī.} — Mira shuddered when she saw the snake.
  \item \emph{Recall:} say the Hindi for to frighten, then the Hindi for cowardly, then say \emph{siharnā} again
\end{itemize}

\begin{tcolorbox}[breakable,colback=teal!4,colframe=teal!35!black,title={Before you move on}]
What does \hi{सिहरना} mean? (\textquotedblleft{}To shudder; to shiver\textquotedblright{}.) Say it once more.
\end{tcolorbox}
